% GENERATED FILE. Edit canonical lessons, then run npm run generate:books.
% canonical-source-hash: fnv1a64:349e677c2e42b2ca
% canonical-lessons: TA-C394-orutalaippatcam, TA-C394-ulnokkam, TA-C394-natunilai, TA-C394-carpaka, TA-C394-kurittu

\chapter{Taking Sides}
\label{ch:ta-taking-sides}
\hlchaptermodality{394}

\begin{chapteropening}
\textbf{By the end of this chapter:} \emph{I can talk about bias, hidden motives and staying neutral.}

Complete the last lesson of chapter 394: I can talk about bias, hidden motives and staying neutral.
\end{chapteropening}

\section[\texorpdfstring{\ta{ஒருதலைப்பட்சம்} (orutalaippaṭcam)}{ஒருதலைப்பட்சம் (orutalaippaṭcam)}]{\ta{ஒருதலைப்பட்சம்} (orutalaippaṭcam) — one-sidedness; bias}

\label{lesson:TA-C394-orutalaippatcam}



\begin{quote}
Before the new one: say the Tamil for roughly, then the Tamil for mostly.
\end{quote}

\subsection*{You'll want to know: \ta{ஒருதலைப்பட்சம்}}
\textbf{\ta{ஒருதலைப்பட்சம்}} — \emph{orutalaippaṭcam} — \textquotedblleft{}one-sidedness; bias\textquotedblright{}.

\begin{quote}
\ta{தீர்ப்பு} \ta{ஒருதலைப்பட்சமாக} \ta{இருந்தால்} \ta{மக்கள்} \ta{ஏற்க} \ta{மாட்டார்கள்}.
\end{quote}

\begin{quote}
\emph{tīrppu orutalaippaṭcamāka iruntāl makkaḷ ēṟka māṭṭārkaḷ.}
\end{quote}

\begin{quote}
If the verdict is one-sided, people will not accept it.
\end{quote}

\ta{ஒருதலைப்பட்சமாக}, one-sidedly: taking only one side.

\begin{cousinweb}
\textbf{\ta{ஒருதலைப்பட்சம்}} (\emph{orutalaippaṭcam}) is built from native \emph{oru}, \textquotedblleft{}one\textquotedblright{} (DEDR 990), and \emph{talai}, \textquotedblleft{}head, end\textquotedblright{} (DEDR 3103), with \emph{paṭcam}, \textquotedblleft{}side\textquotedblright{}, borrowed from Sanskrit \emph{pakṣa} (Monier-Williams): taking one end only.
\end{cousinweb}

\begin{grammarlens}[title={What you've built}]
One more word for saying what happens if.
\end{grammarlens}

\subsection*{Your turn}
\begin{itemize}
\raggedright
  \item \emph{Say it:} \emph{orutalaippaṭcam}
  \item \emph{Say it:} \emph{orutalaippaṭcam}, once more
  \item \emph{Say it:} \emph{tīrppu orutalaippaṭcamāka iruntāl makkaḷ ēṟka māṭṭārkaḷ.} — If the verdict is one-sided, people will not accept it.
  \item \emph{Recall:} say the Tamil for roughly, then the Tamil for mostly, then say \emph{orutalaippaṭcam} again
\end{itemize}

\begin{tcolorbox}[breakable,colback=teal!4,colframe=teal!35!black,title={Before you move on}]
What does \ta{ஒருதலைப்பட்சம்} mean? (\textquotedblleft{}One-sidedness; bias\textquotedblright{}.) Say it once more.
\end{tcolorbox}
\section[\texorpdfstring{\ta{உள்நோக்கம்} (uḷnōkkam)}{உள்நோக்கம் (uḷnōkkam)}]{\ta{உள்நோக்கம்} (uḷnōkkam) — a hidden motive}

\label{lesson:TA-C394-ulnokkam}



\begin{quote}
Before the new one: say the Tamil for mostly, then the Tamil for one-sidedness.
\end{quote}

\subsection*{You'll want to know: \ta{உள்நோக்கம்}}
\textbf{\ta{உள்நோக்கம்}} — \emph{uḷnōkkam} — \textquotedblleft{}a hidden motive\textquotedblright{}.

\begin{quote}
\ta{அவருடைய} \ta{உதவிக்குப்} \ta{பின்னால்} \ta{ஒரு} \ta{உள்நோக்கம்} \ta{இருக்கிறது} \ta{என்று} \ta{நினைக்கிறேன்}.
\end{quote}

\begin{quote}
\emph{avaruṭaiya utavikkup piṉṉāl oru uḷnōkkam irukkiṟatu eṉṟu niṉaikkiṟēṉ.}
\end{quote}

\begin{quote}
I think there is a hidden motive behind his help.
\end{quote}

\ta{உள்}, inner + \ta{நோக்கம்}, aim: an aim kept inside.

\begin{cousinweb}
\textbf{\ta{உள்நோக்கம்}} (\emph{uḷnōkkam}) is native: \emph{uḷ}, \textquotedblleft{}inside, mind\textquotedblright{} (DEDR 698), and \emph{nōkkam}, \textquotedblleft{}aim, purpose\textquotedblright{}, from \emph{nōkku}, \textquotedblleft{}to look\textquotedblright{} (DEDR 3794): an aim the eye does not see.
\end{cousinweb}

\begin{grammarlens}[title={What you've built}]
One more word for saying what happens if.
\end{grammarlens}

\subsection*{Your turn}
\begin{itemize}
\raggedright
  \item \emph{Say it:} \emph{uḷnōkkam}
  \item \emph{Say it:} \emph{uḷnōkkam}, once more
  \item \emph{Say it:} \emph{avaruṭaiya utavikkup piṉṉāl oru uḷnōkkam irukkiṟatu eṉṟu niṉaikkiṟēṉ.} — I think there is a hidden motive behind his help.
  \item \emph{Recall:} say the Tamil for mostly, then the Tamil for one-sidedness, then say \emph{uḷnōkkam} again
\end{itemize}

\begin{tcolorbox}[breakable,colback=teal!4,colframe=teal!35!black,title={Before you move on}]
What does \ta{உள்நோக்கம்} mean? (\textquotedblleft{}A hidden motive\textquotedblright{}.) Say it once more.
\end{tcolorbox}
\section[\texorpdfstring{\ta{நடுநிலை} (naṭunilai)}{நடுநிலை (naṭunilai)}]{\ta{நடுநிலை} (naṭunilai) — neutrality; neutral}

\label{lesson:TA-C394-natunilai}



\begin{quote}
Before the new one: say the Tamil for one-sidedness, then the Tamil for a hidden motive.
\end{quote}

\subsection*{You'll want to know: \ta{நடுநிலை}}
\textbf{\ta{நடுநிலை}} — \emph{naṭunilai} — \textquotedblleft{}neutrality; neutral\textquotedblright{}.

\begin{quote}
\ta{நீதிபதி} \ta{எப்போதும்} \ta{நடுநிலையாக} \ta{இருக்க} \ta{வேண்டும்}.
\end{quote}

\begin{quote}
\emph{nītipati eppōtum naṭunilaiyāka irukka vēṇṭum.}
\end{quote}

\begin{quote}
A judge must always be neutral.
\end{quote}

\ta{நடுநிலையாக}, impartially.

\begin{cousinweb}
\textbf{\ta{நடுநிலை}} (\emph{naṭunilai}) is native: \emph{naṭu}, \textquotedblleft{}middle, centre, equity\textquotedblright{} (DEDR 3584), and \emph{nilai}, \textquotedblleft{}a standing\textquotedblright{} (DEDR 3675): standing in the middle.
\end{cousinweb}

\begin{grammarlens}[title={What you've built}]
One more word for saying what happens if.
\end{grammarlens}

\subsection*{Your turn}
\begin{itemize}
\raggedright
  \item \emph{Say it:} \emph{naṭunilai}
  \item \emph{Say it:} \emph{naṭunilai}, once more
  \item \emph{Say it:} \emph{nītipati eppōtum naṭunilaiyāka irukka vēṇṭum.} — A judge must always be neutral.
  \item \emph{Recall:} say the Tamil for one-sidedness, then the Tamil for a hidden motive, then say \emph{naṭunilai} again
\end{itemize}

\begin{tcolorbox}[breakable,colback=teal!4,colframe=teal!35!black,title={Before you move on}]
What does \ta{நடுநிலை} mean? (\textquotedblleft{}Neutrality; neutral\textquotedblright{}.) Say it once more.
\end{tcolorbox}
\section[\texorpdfstring{\ta{சார்பாக} (cārpāka)}{சார்பாக (cārpāka)}]{\ta{சார்பாக} (cārpāka) — on behalf of; in favour of}

\label{lesson:TA-C394-carpaka}



\begin{quote}
Before the new one: say the Tamil for a hidden motive, then the Tamil for neutrality.
\end{quote}

\subsection*{You'll want to know: \ta{சார்பாக}}
\textbf{\ta{சார்பாக}} — \emph{cārpāka} — \textquotedblleft{}on behalf of; in favour of\textquotedblright{}.

\begin{quote}
\ta{என்} \ta{சார்பாக} \ta{என்} \ta{அண்ணன்} \ta{கூட்டத்துக்குப்} \ta{போனான்}.
\end{quote}

\begin{quote}
\emph{eṉ cārpāka eṉ aṇṇaṉ kūṭṭattukkup pōṉāṉ.}
\end{quote}

\begin{quote}
My elder brother went to the meeting on my behalf.
\end{quote}

X-\ta{இன்} \ta{சார்பாக}, on behalf of X; X-\ta{க்குச்} \ta{சார்பாக}, in favour of X.

\begin{cousinweb}
\textbf{\ta{சார்பாக}} (\emph{cārpāka}) is native: \emph{cārpu}, \textquotedblleft{}support, side\textquotedblright{}, from \emph{cār}, \textquotedblleft{}to lean on, depend upon, take shelter in\textquotedblright{} (DEDR 2460), with the adverb ending -\ta{ஆக}.
\end{cousinweb}

\begin{grammarlens}[title={What you've built}]
One more word for saying what happens if.
\end{grammarlens}

\subsection*{Your turn}
\begin{itemize}
\raggedright
  \item \emph{Say it:} \emph{cārpāka}
  \item \emph{Say it:} \emph{cārpāka}, once more
  \item \emph{Say it:} \emph{eṉ cārpāka eṉ aṇṇaṉ kūṭṭattukkup pōṉāṉ.} — My elder brother went to the meeting on my behalf.
  \item \emph{Recall:} say the Tamil for a hidden motive, then the Tamil for neutrality, then say \emph{cārpāka} again
\end{itemize}

\begin{tcolorbox}[breakable,colback=teal!4,colframe=teal!35!black,title={Before you move on}]
What does \ta{சார்பாக} mean? (\textquotedblleft{}On behalf of; in favour of\textquotedblright{}.) Say it once more.
\end{tcolorbox}
\section[\texorpdfstring{\ta{குறித்து} (kuṟittu)}{குறித்து (kuṟittu)}]{\ta{குறித்து} (kuṟittu) — about; concerning}

\label{lesson:TA-C394-kurittu}



\begin{quote}
Before the new one: say the Tamil for neutrality, then the Tamil for on behalf of.
\end{quote}

\subsection*{You'll want to know: \ta{குறித்து}}
\textbf{\ta{குறித்து}} — \emph{kuṟittu} — \textquotedblleft{}about; concerning\textquotedblright{}.

\begin{quote}
\ta{இது} \ta{குறித்துக்} \ta{கேள்வி} \ta{இருந்தால்} \ta{ஆசிரியரிடம்} \ta{கேட்கலாம்}.
\end{quote}

\begin{quote}
\emph{itu kuṟittuk kēḷvi iruntāl āciriyariṭam kēṭkalām.}
\end{quote}

\begin{quote}
If you have questions about this, you can ask the teacher.
\end{quote}

A written word for \ta{பற்றி}: X \ta{குறித்து}, about X.

\begin{cousinweb}
\textbf{\ta{குறித்து}} (\emph{kuṟittu}) is native: the verbal participle of \emph{kuṟi}, \textquotedblleft{}to aim at, intend, point to\textquotedblright{} (DEDR 1847): \textquotedblleft{}having aimed at\textquotedblright{}, so \textquotedblleft{}concerning\textquotedblright{}.
\end{cousinweb}

\begin{grammarlens}[title={What you've built}]
One more word for saying what happens if.
\end{grammarlens}

\subsection*{Your turn}
\begin{itemize}
\raggedright
  \item \emph{Say it:} \emph{kuṟittu}
  \item \emph{Say it:} \emph{kuṟittu}, once more
  \item \emph{Say it:} \emph{itu kuṟittuk kēḷvi iruntāl āciriyariṭam kēṭkalām.} — If you have questions about this, you can ask the teacher.
  \item \emph{Recall:} say the Tamil for neutrality, then the Tamil for on behalf of, then say \emph{kuṟittu} again
\end{itemize}

\begin{tcolorbox}[breakable,colback=teal!4,colframe=teal!35!black,title={Before you move on}]
What does \ta{குறித்து} mean? (\textquotedblleft{}About; concerning\textquotedblright{}.) Say it once more.
\end{tcolorbox}
